\section{问题二模型的建立与求解}

\subsection{问题二的描述与模型继承}

问题二在问题一的基础上增加两个条件：长城站外风速为$3\ \mathrm{m/s}$，实验者由静止站立改为轻微运动。三层防护材料的厚度、密度、比热、导热系数以及功能层温度相关放热函数均未改变，因此问题二继续采用问题一的人体单节点热容--三层一维瞬态导热耦合模型，只修改人体侧和服装外侧两个对流边界。

附件将人体状态与人体--内层织物之间的换热系数直接对应：静止站立取$3.0\ \mathrm{W/(m^2\cdot K)}$，轻微运动取$4.0\ \mathrm{W/(m^2\cdot K)}$。因此本问以附件给出的状态修正作为轻微运动对人体侧换热的直接表征，不再额外引入题目未给定的代谢产热参数。

\subsection{有风条件下的外表面对流修正}

问题一无外加风速，附件给出的
\begin{equation}
h_{\rm nat}=2.38\left|T_{\rm cl}-T_\infty\right|^{0.25}
\end{equation}
对应服装表面--环境之间的自然对流支路，其中$T_{\rm cl}=T_3(L,t)$。当外部存在有限风速时，服装表面对流还受到强制对流控制。ISO~7730中服装表面对流换热系数采用自然对流与强制对流两支关系中的较大者\cite{iso7730}：
\begin{equation}
h_e(t)=\max\left\{
2.38\left|T_3(L,t)-T_\infty\right|^{0.25},
\ 12.1\sqrt{v}
\right\}.
\label{eq:q2-hout-mixed}
\end{equation}
该写法保留了温差对自然对流的作用；在强制对流占优时，温差仍通过对流热流
\begin{equation}
q_e''=h_e\left[T_3(L,t)-T_\infty\right]
\end{equation}
进入能量交换，并未从传热模型中消失。

近年的热人模研究表明，全身对流换热系数随空气速度按幂函数显著增加，且$3\ \mathrm{m/s}$处于相关实验/数值研究的有效风速区间内\cite{yang2023hc}。穿衣人体的CFD研究还表明，外部风会通过服装开口和织物附近流场增强衣下热交换，因此有风条件下不能继续沿用静止空气的单一自然对流关系\cite{xu2022clothingwind}。

本问$v=3\ \mathrm{m/s}$，强制对流支路为
\begin{equation}
h_{\rm for}=12.1\sqrt{3}=20.96\ \mathrm{W/(m^2\cdot K)}.
\label{eq:q2-hforced}
\end{equation}
实验者进入室外时服装系统初始温度为$37\,^{\circ}\mathrm C$，环境温度为$T_\infty=-40\,^{\circ}\mathrm C$，故初始最大温差为$77\ \mathrm K$，此时自然对流支路仅为
\begin{equation}
h_{\rm nat,max}=2.38\times77^{0.25}\approx7.05\ \mathrm{W/(m^2\cdot K)}.
\end{equation}
随后服装表面温度下降，自然对流支路进一步减小。因此在本问的整个主要降温过程中均有$h_{\rm for}>h_{\rm nat}$，式\eqref{eq:q2-hout-mixed}退化为
\begin{equation}
h_e=20.96\ \mathrm{W/(m^2\cdot K)}.
\label{eq:q2-hout-value}
\end{equation}
于是服装外表面的Robin边界为
\begin{equation}
-k_3T_{3,x}(L,t)
=20.96\left[T_3(L,t)-T_\infty\right]
=20.96\left[T_3(L,t)+40\right].
\label{eq:q2-outer-robin}
\end{equation}

\subsection{轻微运动条件下的人体侧修正}

附件给出轻微运动时人体表面--内层织物之间的对流换热系数
\begin{equation}
h_b^{(2)}=4.0\ \mathrm{W/(m^2\cdot K)}.
\label{eq:q2-hbody}
\end{equation}
因此人体侧Robin边界由问题一的$h_b=3.0$改为
\begin{equation}
-k_1T_{1,x}(0,t)
=4.0\left[T_h(t)-T_1(0,t)\right].
\label{eq:q2-inner-robin}
\end{equation}
人体等效热容$C_h$和有效面积$A_b$沿用问题一，人体能量平衡改写为
\begin{equation}
C_h\frac{\mathrm dT_h}{\mathrm dt}
=-4.0A_b\left[T_h(t)-T_1(0,t)\right],
\qquad T_h(0)=37\,^{\circ}\mathrm C.
\label{eq:q2-body-ode}
\end{equation}

\subsection{问题二完整传热模型}

三层材料内部控制方程与问题一相同：
\begin{equation}
\left\{
\begin{aligned}
\rho_1c_1T_{1,t}&=k_1T_{1,xx},\\
\rho_2c_2T_{2,t}&=k_2T_{2,xx}+\rho_2q_{\rm pcm}(T_2),\\
\rho_3c_3T_{3,t}&=k_3T_{3,xx}.
\end{aligned}
\right.
\label{eq:q2-governing}
\end{equation}
初始条件仍为
\begin{equation}
T_h(0)=37\,^{\circ}\mathrm C,
\qquad
T_i(x,0)=37\,^{\circ}\mathrm C,\quad i=1,2,3.
\end{equation}
两个材料界面继续满足
\begin{equation}
\left\{
\begin{aligned}
&T_1(x_1^-,t)=T_2(x_1^+,t),\\
&k_1T_{1,x}(x_1^-,t)=k_2T_{2,x}(x_1^+,t),\\
&T_2(x_2^-,t)=T_3(x_2^+,t),\\
&k_2T_{2,x}(x_2^-,t)=k_3T_{3,x}(x_2^+,t).
\end{aligned}
\right.
\end{equation}
再与式\eqref{eq:q2-inner-robin}、\eqref{eq:q2-body-ode}和\eqref{eq:q2-outer-robin}联立，即得到问题二闭合模型。

\subsection{求解方法的继承与简化}

问题二继续采用问题一的微小时间段拓展分离变量法。功能层$T_2$的Dirichlet--Dirichlet展开和$q_{\rm pcm}(T)$的时间步更新完全保留；第一层仍为Robin--Dirichlet问题，只需将本征值方程中的$h_b$由$3.0$替换为$4.0$：
\begin{equation}
\tan(\lambda_{1,m}d_1)=-\frac{k_1\lambda_{1,m}}{4.0}.
\end{equation}
第三层仍为Dirichlet--Robin问题，但由于$3\ \mathrm{m/s}$下强制对流支路始终占优，$h_e=20.96\ \mathrm{W/(m^2\cdot K)}$可在整个计算过程中保持不变，因此第三层本征值也只需求解一次：
\begin{equation}
\tan(\lambda_{3,m}d_3)=-\frac{k_3\lambda_{3,m}}{20.96}.
\end{equation}
相较问题一，问题二不再需要逐时间步更新外侧自然对流系数，其余时间递推、界面热流连续和人体能量平衡的求解流程保持一致。

\subsection{返回室内时间}

问题二仍以人体等效表面温度首次降至$15\,^{\circ}\mathrm C$作为必须返回室内的时刻：
\begin{equation}
t_{15}^{(2)}=\inf\{t>0:T_h(t)\le15\,^{\circ}\mathrm C\}.
\end{equation}
同时记录
\begin{equation}
t_{10}^{(2)}=\inf\{t>0:T_h(t)\le10\,^{\circ}\mathrm C\}
\end{equation}
作为危险暴露时间。数值结果在完成时间步和模态截断收敛检查后写入正文。
